
\subsection{Measurement grid}
\label{sec:scaling-grid}

We measure throughput of the multi-query \texttt{pmap} configuration
from Section~\ref{sec:parallelism} on a full grid of chip count
(1, 2, 4, 8) and sequence length (100, 250, 500, 1000 residues), 16
points total, each a single steady-state measurement with no
repeat.
At chips=8, throughput falls from 19.153 to 0.522 proteins/second as
length grows from 100 to 1000 residues. At length=100, it rises from
2.935 to 19.153 proteins/second as chips grow from 1 to 8. Every point
on this grid uses the rotated-alphabet input described in
Section~\ref{sec:methodology}, not the 118-residue toy sequence used
for the single-query experiments elsewhere in the paper. Because chip
count and length are varied on the same input family and the same
timer, this grid is the paper's cleanest measure of what adding chips
buys.

\begin{figure}[H]
  \centering
  \includegraphics[width=\linewidth]{scaling_grid.pdf}
  \caption{The 16-point scaling grid (multi-query \texttt{pmap},
  repeated-alphabet inputs, one process per point). (a) Measured
  throughput against chips in use for each sequence length, with the
  fitted law $4527.77 \times \mathrm{chips}^{0.963} \times
  \mathrm{length}^{-1.572}$ drawn dashed. (b) Fitted divided by measured
  throughput at every grid point: the fit over-predicts at 100 and 1000
  residues and under-predicts at 250 and 500, a systematic pattern. The one-chip point at 1000 residues is recorded
  to two significant figures.}
  \label{fig:scaling-grid}
\end{figure}

\subsection{Measured speedup at each length}
\label{sec:scaling-speedup}

The eight-over-one-chip speedup
is 6.53$\times$ at length 100, rising to 7.47$\times$ at 250,
7.76$\times$ at 500, and 7.91$\times$ at 1000.
The four-over-one and two-over-one speedups show the same pattern:
3.55$\times$ to 3.97$\times$, and 1.88$\times$ to 1.98$\times$,
respectively, moving toward their ideal values (4$\times$, 2$\times$)
as length grows.
Scaling efficiency is worst at the shortest sequence length in this
grid, which is also closest to the 118-residue input used
throughout the rest of the paper: the fixed per-call overhead that
Sections~\ref{sec:profiling} and~\ref{sec:parallelism} describe
makes up a larger share of a short call, leaving less of the
eight-chip speedup to be realized. At longer sequences, where
per-chip compute dominates more of the call, the measured speedup
approaches ideal linear scaling.

\subsection{A descriptive power-law fit}
\label{sec:scaling-fit}

An ordinary-least-squares fit in log-log space summarizes the grid as
\begin{equation}
  \text{throughput} \approx 4527.77 \times \text{chips}^{0.963} \times \text{length}^{-1.572},
  \quad R^2 = 0.981.
  \label{eq:scaling-fit}
\end{equation}
The chip exponent, 0.963, is close enough to 1 to describe scaling
across chips as near-linear over this range. The length exponent,
$-1.572$, is steeper than $-1$: throughput falls faster than simple
inverse proportionality with sequence length, consistent with the
pair representation's cost growing faster than linearly in residue
count (Section~\ref{sec:background}).

$R^2 = 0.981$ describes the fit in log space across all 16 points; it
does not mean the fit tracks the grid closely at any one length.
The ratio of fitted to measured throughput follows a systematic
pattern: the fit over-predicts
by 10--26\% at length 100 and by 23--32\% at length 1000, and
under-predicts by 11--13\% at length 250 and by 22--26\% at length
500 (Figure~\ref{fig:scaling-grid}).
A single global exponent therefore approximates a relationship that
curves, and we use the fit only as a summary of this grid. We give
per-length residuals instead of a regression interval on the
exponents: with one measurement per cell the deviation is dominated by
misspecification, not by sampling noise, and an interval would
understate it. A chip exponent close to 1 predicts that cost per
prediction is roughly flat in chip count. Section~\ref{sec:cost} tests
this on the measured grid and finds it holds only at longer
sequences.

